\section{The continuation and its exact scope}
\label{sec:scope}

\subsection{What problem is being improved?}

The long-term objective is a complete algorithm deciding whether a classical
knot diagram with $n$ crossings represents the unknot in time
$n^{O(\log n)}=2^{O((\log n)^2)}$. This is a statement about the entire
computation, including preprocessing, state construction, every restart,
coefficient bit lengths, and the final verdict. More generally,
quasi-polynomial means $\exp((\log(n+2))^{O(1)})$; we focus on the
concrete squared-logarithm target. A polynomial subroutine
for an exponentially large state is not sufficient.

The maintained implementation combines inexpensive certificates and
invariant filters with complete, potentially exponential fallback
algorithms. Its Khovanov route uses a tangles-and-cobordisms formulation
of chain complexes, following the framework of Bar-Natan~\cite{barnatan};
the unknot-detection theorem itself is due to Kronheimer and
Mrowka~\cite{km}. Exactness and the worst-case running time are separate
questions. In particular, a filter which has not detected an obstruction
has not proved the knot trivial.

This continuation audits the repository at \sourcepin, fetched on
8 October 2026. The initial remote inspection saw the immediately earlier
revision \code{04d799b69d38}; the branch advanced
while the sparse checkout was being obtained. All patches, baseline
comparisons, and distributed source hashes use the actual checked-out
revision \sourcepin. No inference is made from a mutable \code{main} URL
after that point. The full baseline identifier is:
\begin{center}
\code{ea9cf443c8ce33e09ed0c48b0ef7856ddf6326a3}
\end{center}

\subsection{The three bottlenecks addressed}

Recent incoming reports already supplied polynomial projectors, compressed
surface-cover assembly, regular-cover gluing equivalence, and a binary
normal-disc certificate checker~\cite{projectors,certificates}. We use
their precise scope, including their negative measurements. The present
work advances three subsequent tasks.

\begin{enumerate}
\item \textbf{Optimize over a represented family.} Previous surface-cover
kernels primarily classify a supplied assignment or compare two assignments.
We ask an existential and optimization question over all assignments
satisfying simultaneous linear congruences. This removes a genuine phase
enumeration for the specified connectivity predicate.
\item \textbf{Extract compressed data from actual normal coordinates.}
Rather than assume an already supplied sheet presentation, we construct
signed disc-stack and prism-gap identifications from tetrahedral face
pairings and seven normal coordinates per tetrahedron. This gives a
concrete geometric interface, while stopping short of full compressed
cutting and parallelity-bundle classification.
\item \textbf{Reuse algebra already computed.} The projector report
identified the corner algebra $e\End(C)e$ as the correct child
endomorphism space but did not implement reuse. We implement and measure
that inheritance, including canonicalization needed to preserve the
existing bounded candidate policy.
\end{enumerate}

\begin{center}
\small
\begin{tabular}{>{\raggedright\arraybackslash}p{.25\linewidth}>{\raggedright\arraybackslash}p{.34\linewidth}>{\raggedright\arraybackslash}p{.30\linewidth}}
\toprule
Contribution & What is established & Remaining limitation\\
\midrule
Affine cyclic families & Exact minimum component count; factor-free
construction; dual certificates & Translation actions and affine
constraints; no arbitrary topology predicate\\
Normal-coordinate extraction & Linear number of signed bands and local
prism contacts; exact binary arithmetic & No full residual handle
structure or unrestricted orbit implementation\\
Scalar-space inheritance & Exact child spaces from one parent solve;
verified implementation & Parent dimension, component size, and splitter
activation still matter\\
\bottomrule
\end{tabular}
\end{center}

\subsection{Relation to the current geometric literature}

Lackenby's July 2026 preprint~\cite{lackenby2026} supplies an
incompressibility algorithm and an unknot-recognition consequence. Its
Section 9 does not provide cost bounds for every step; Proposition 9.1
instead bounds iterations by $L(g+1)^L$. Theorem 14.4 already gives a
polynomial-time compressed cutting operation for a uniform-type handle
structure and a specified orientable normal surface, measured in handle
count and logarithmic weight. Thus neither compressed cutting in principle
nor the use of binary surface coordinates is a new theorem of this
article. The new extraction code makes a smaller, explicit tetrahedral
part of that interface executable.

Agol--Hass--Thurston's interval-orbit methods~\cite{aht} are the appropriate
general background for processing enormous collections of normal pieces.
Our local bounds are not a claimed improvement on their general orbit
complexity. Covering-space monodromy and its connection with connected
components are classical~\cite{hatcher}. The affine-family theorem adds
a fully specified optimization and certificate layer for a restricted
family representation.

\subsection{Meaning of ``proved'' and of the reported novelty}

Theorems below have conventional mathematical proofs. The arithmetic,
geometric, and algebraic implementations are supported by exact tests and
independent finite oracles. Those are different forms of evidence: a
finite test does not prove an asymptotic theorem, and a prose proof does
not verify a Python implementation. Separate model-assisted reviews were
used to challenge the proofs and check implementation contracts. No
proof-assistant formalization or external peer review is claimed.

The delta relative to the inspected repository and incoming reports is
explicit. We do not claim worldwide priority for the underlying facts
about affine modules, primitive vectors, finite-ring duality, or corner
algebras. Elementary ingredients can be useful even when they are
classical; the research contribution here is the precise executable
combination, its certificates, its performance consequences, and the
identification of what it does and does not resolve.
